\documentclass[10pt]{article}

\usepackage[utf8]{inputenc}

\usepackage[T1]{fontenc}

\usepackage{amsmath}

\usepackage{amsfonts}

\usepackage{amssymb}

\usepackage[version=4]{mhchem}

\usepackage{stmaryrd}

\usepackage{bbold}



\begin{document}

скопич. состояние системы в данный момент времени полностью определяется нек-рой функцией распределения $f(q, p)$. Наблюдаемое макроскопич. значение $A(\vec{x}, t)$ микроскопич. динамич. переменной $a(p, q, \vec{x}, t)$, по определению, выражается соотношением (4). Если $f(q, p) \geqslant 0$, то Р. ф. можно интерпретировать как плотность вероятности нахождения системы в точке фазового пространства $(q, p)$. Конкретный выбор функции $f(q, p)$ не является задачей механики. Она должна наиболее полно отображать доступную наблюдателю информацию о состоянии системы. На практике при исследовании системы $N$ тождественных частиц нормировка (5) не всегда удобна, так как не обеспечивает корректный предельный переход от квантовой статистики к классической. Поэтому предпочтительнее использовать юрмировку



\[

\int f(q ; p) d \Gamma=1,

\]



где $d \Gamma=\frac{d q d p}{N !(2 \pi \hbar)^{3 N}}-$ безразмерный элемент фазового объема. Тогда макроскопич. значение любой микроскопич. динамич. переменной будет определяться соотношением



\[

A(\vec{x}, t)=\int d \Gamma a(p, q, \vec{x}, t) f(q, p) .

\]



Интегрирование в (6) соответствует суммированию по всем различным состояниям. Тем самым учитывается, что перестановка тождественных частиц в квантовой механике не меняет состояния и это свойство должно сохраняться и в предельном случае классич. статистич. механики. Множитель $1 / n !$ при фазовом объеме был введен Дж. Гиббсом (J. Gibbs), что позволило разрешить одноименный парадокс, заключающийся в возрастании энтропии при смешении одинаковых газов при одинаковых температуре и давлении.



A. В. Солдатов.



РАСПРЕДЕЛЕНИЯ ФУНКЦИЯ ДлЯ статистиче ских с истем - плотность распределения вероятностей (плотность вероятностной меры относительно меры Лебега), заданного на фазовом пространстве классической статистической системы частиц. Иными словами,



\[

f_{N}\left(t, x_{1}, p_{1} ; \ldots ; x_{N}, p_{N}\right) d \Gamma_{N}

\]



где $x_{i}$ и $p_{i} \in \mathbb{R}^{v}, v=1,2,3,-$ соответственно координата и импульс $i$-й частицы, $d \Gamma_{N}=d x_{1} \ldots d x_{N} d p_{1} \ldots d p_{N}$, а $N-$ число частиц,- это вероятность обнаружить в момент времени $t 1-ю, \quad \ldots, \quad N$-ю частицу в малых объемах $d x_{1}, d p_{1}, \ldots, d x_{N}, d p_{N}$ с центрами, соответственно, в точках $x_{1}, p_{1}, \ldots, x_{N}, p_{N}$. Для системы тождественных частиц в рамках канонического ансамбля Р. Ф. нормированы следующим образом:



\[

\frac{1}{\hbar^{3 N} N !} \int d \Gamma_{N} f_{N}\left(t, x_{1}, \ldots, x_{N}, p_{1}, \ldots, p_{N}\right)=1

\]



( $\hbar$ - постоянная Планка).



Последовательность Р. Ф. большого канонического ансамбля удовлетворяет следующему условию нормировки:



\[

f_{0}+\sum_{N \geqslant 1} \frac{1}{\hbar^{3 N} N !} \int d \Gamma_{N} f_{N}\left(t, x_{1}, p_{1} ; \ldots ; x_{N}, p_{N}\right)=1 .

\]



Удобным способом описания состояний статистич. систем, то есть задания соответствуюшего распределения вероятностей на фазовом пространстве системы, служат также последовательности распределения вероятностей частичных функций (иначе: приведенных функций распреде ле н и я), удовлетворяющие согласованности условиям. В канонич. ансамбле приведенные функции распределения определяются формулой



\[

\begin{gathered}

F_{s}^{(N)}\left(x_{1}, p_{1} ; \ldots ; x_{s}, p_{s}\right)= \\

=\frac{N !}{(N-s) !} \int d \Gamma_{N-s} f_{N}\left(t, x_{1}, p_{1} ; \ldots ; x_{N}, p_{N}\right), \quad 1 \leqslant s \leqslant N,

\end{gathered}

\]



а в большом канонич. ансамбле - формулой



\[

F_{s}^{(N)}\left(t, x_{1}, p_{1} ; \ldots ; x_{s}, p_{s}\right)=

\]



\[

=1+\sum_{N=1}^{\infty} \frac{1}{N !} \int d x_{s+1} d p_{s+1} \cdots

\][



··· d x\_\{s+N\} d p\_\{s+N\} f\_\{N\}\textbackslash left(t, x\_\{1\}, p\_\{1\} ; ··· ; x\_\{N\}, p\_\{N\}\textbackslash right), \textbackslash quad s \textbackslash geqslant 1 .

]



Частичные Р. Ф. соответствуют моментным функциям $s$-го порядка для случайных полей. Частичные функции распределения называются также корреляцион ны м и функц и я м и. В квантовой механике роль Р. Ф. и частичных функций распределения играют, соответственно, плотности матрицы и приведенные матрицы плотности.



Лит. см. при ст. Гиббсовские состояния. РАСПРЕДЕЛЕНИЯ ФУНКЦИЯ с Л У ч а й н о й в е л и ч и н ы $X$ - функция действительного переменного $x$, принимающая при каждом $x$ значение, равное вероятности неравенства $X<x$.



Каждая Р. Ф. $F(x)$ обладает следующими свойствами:



\begin{enumerate}

  \item $F\left(x^{\prime}\right) \leqslant F\left(x^{\prime \prime}\right)$ при $x^{\prime}<x^{\prime \prime}$;



  \item $F(x)$ непрерывна слева при каждом $x$;



  \item $\lim _{x \rightarrow-\infty} F(x)=0, \quad \lim _{x \rightarrow \infty} F(x)=1$



\end{enumerate}



(иногда Р. ф. определяют как вероятность неравенства $X \leqslant x$, и тогда она оказывается непрерывной справа). Каждая функция $F$, обладающая свойствами 1)-3), может рассматриваться как Р. Ф. нек-рой случайной величины $X$.



В. И. Битюцков.



PACПPOCTPAHEHИE BOЛН в турбуле н т ны с с е да х сопровождается преломлением волн на турбулентных неоднородностях показателя преломления $n^{\prime}$ (для распространения электромагнитных волн в атмосфере, создаваемых турбулентными флуктуациями давления, температуры и влажности воздуха), а также флуктуациями скорости распространения волн (для звуковых волн в атмосфере, создаваемых, главным образом, флуктуациями температуры). В результате возникает рассеяние волн, а также пульсации амплитуды и фазы волн, прошедших через среду, что создает мерцание и дрожание изображений источников излучения в приемных устройствах (могущее, к тому же, быть различным в разных участках длин волн, то есть, в случае видимого света, к пульсациям цвета источника). Эти эффекты создают помехи для оптич. астрономии, радиоастрономии, оптической и радиосвязи с искусственными спутниками Земли и космич. ракетами и для гидроакустич. связи в море.



Эффективное сечение рассеяния электромагнитных волн в атмосфере объемом $V$ в телесный угол $d \Omega$ около направления $\boldsymbol{q}$ получается в виде:



\[

d \sigma=\left(4 \pi^{2}\right)^{-1} k^{4} \sin ^{2} \alpha G G^{*} d \Omega,

\]



где $G(\boldsymbol{k}-k \boldsymbol{q})$ - преобразование Фурье флуктуаций коэффициента преломления, $\boldsymbol{k}-$ волновой вектор падающей волны, а $\alpha-$ угол между $q$ и вектором поляризации падающей волны. Осреднение дает



\[

\overline{G G^{*}}=8 \pi^{3} V F(\boldsymbol{k}-k \boldsymbol{q}),

\]



где $F$ - спектральная функция показателя преломления. Таким образом, в рассеянии на данный угол участвуют турбулентные неоднородности, образующие пространственную дифракционную решетку с периодом $2 \pi|\boldsymbol{k}-k \boldsymbol{q}|^{-1}$. В инерционном интервале спектра турбулентности получается



\[

\overline{d \sigma} \sim\left(\sin \frac{\theta}{2}\right)^{-11 / 3}

\]



где $\theta$ - угол рассеяния.



Для расчета флуктуаций амплитуды и фазы волны используют уравнение для эйконала $\psi$ в приближении «плавных возмущений»



\[

\Delta \psi^{\prime}+2 i k \nabla \psi^{\prime}=-2 k^{2} n^{\prime}

\]



из к-рого следует, что флуктуации эйконала получаются интегрированием $n^{\prime}$ с нек-рыми весами по рассеивающему объему. Приближенное вычисление соответствующих интегралов осуществляется по-разному в зависимости от ве-





\end{document}